\rulesection{15.0 Entrenchments}

Combat units may entrench to gain defensive benefits.

\rulesubsection{15.1 Entrenchments}

Entrenchment requires two full friendly Movement Phases. During each of those phases the unit must not expend MPs and must have at least 1 CP. At the end of the first Movement Phase, place a Constructing Entrenchment marker. At the end of the second, flip it to its Entrenchment side.

The two phases must occur on consecutive turns. If a unit with a Constructing Entrenchment marker fails to meet the conditions on the following turn, remove the marker.

\textbf{15.11} Each unit must entrench for itself. A unit entering a hex with an entrenched unit does not also become entrenched. One unit cannot take over another unit's Constructing Entrenchment marker. If an entrenched unit moves, remove its Entrenchment marker.

\textbf{15.12} A unit that attacks remains entrenched if it still occupies the hex after combat, but gains no benefit from entrenchments during that battle.

\rulesubsection{15.2 Effects}

Defending units in entrenchments are doubled when defending and receive a modifier to demoralization die rolls. Cavalry in entrenchments have no ZOC. In addition, the attacker shifts one column left on the Combat Results Table if the defender is entrenched.
